\subsection{Protocol and measurement}
Structured streams recombine public bAbI facts \citep{babi} with independently seeded observation and query schedules. There are four stable facts and two groups of four revisable facts; one group is currently demanded. Observations, queries, and feedback interleave. Context switches and observed revisions are explicit. These are controlled management tasks, not the official bAbI reasoning benchmark. Content templates recur across worlds and earlier development data; new schedule seeds do not establish content-domain generalization.

Methods share input tapes, initial records, slot limits, query budgets, and one overflow decision per eligible observation. Feedback arrives after the answer; the same corruption tape is used, but delivered correctness can differ because methods answer differently. Comparators are fixed exponential recency, retrieval frequency, LRU, frozen-state scoring, a calibrated replacement control, an adaptive rule, and \ame. Frozen state retains learned insertion values but disables subsequent per-item changes. Fixed decay orders by most recent observation, whereas LRU also tracks retrieval. The replacement control adds a seeded admission gate; positive rescaling of decay alone cannot change the eviction order.

Every formal condition uses 24 schedule worlds and three policy/checkpoint seeds (11, 23, 42). Seeds are averaged within a world before inference. We report paired 95\% percentile-bootstrap intervals (5,000 resamples) and sign-flip tests (10,000 draws), with Holm correction over each batch's declared comparisons. Repeated policy seeds are not independent worlds. C1--C4 comprise 17,568 policy trajectories; B1-A and B2--B7 add 5,184 and 7,560, respectively. Each batch has separate mini replay and frozen inputs. Protocol amendments and their preserved superseded evidence are disclosed in Appendix~\ref{app:protocol}.

Accuracy counts correct answers over all scored queries; stale error counts obsolete answers over the same denominator. Stable retention (SR) restricts correctness to stable queries. Adaptation lag is the first position at which two consecutive overlapping four-query windows are each at least 75\% correct; unrecovered episodes are capped at their target-query horizon $H$. Adaptation score is $\mathrm{AS}=1-\mathrm{lag}/H$. It is a bounded first-passage summary, not proof of sustained recovery. Different episode lengths are never pooled without qualification.

\subsection{RQ1: When does state evolution help?}
\begin{table}[t]
\centering\small
\caption{Constrained structured memory: eight slots, sixteen interference keys. Percentages except lag (target queries). All seven policies share the same streams; 24 worlds, three policy seeds. Full intervals accompany the machine-readable evidence (Appendix~\ref{app:tables}).}\label{tab:main}
\begin{tabular}{lrrrr}\toprule
Method & Acc.$\uparrow$ & Stale$\downarrow$ & Ret.$\uparrow$ & Lag$\downarrow$\\\midrule
Fixed decay & 46.48 & 0.00 & 46.64 & 12.88\\
Frequency & 48.50 & 14.03 & 69.76 & 117.17\\
LRU & 41.91 & 0.00 & 42.10 & 17.29\\
Frozen state & 75.46 & 3.96 & 78.84 & 29.72\\
Replacement control$^{\dagger}$ & 54.93 & 1.81 & 74.84 & 103.77\\
Adaptive rule & 71.02 & 0.01 & 79.09 & 42.33\\
AME & 78.06 & 0.02 & 87.48 & 23.03\\
\bottomrule\end{tabular}
\par\vspace{2pt}\footnotesize $\dagger$ Calibrated admission/lifetime control; distributional matching failed. This row is not an exactly matched counterfactual.
\end{table}

With eight slots, \ame improves overall accuracy and stable retention, but fixed decay adapts faster and already achieves zero stale error (Table~\ref{tab:main}). The capacity sweep changes both content slots and interference keys ($2K$), while keeping twelve core identities; it therefore measures \emph{joint capacity/workload scaling}, not a pure scaling law. At four slots, \ame reaches 45.26\% accuracy versus fixed decay's 20.23\%. At sixteen slots, LRU reaches 98.81\%, exceeding \ame's 92.90\%; at 64 slots, their accuracies are 98.99\% and 98.86\%. The advantage is strongest under constrained storage, not across every budget.
Figure~\ref{fig:capacity} shows all budgets. Fixed decay's zero stale error at eight slots coexists with poor stable retention; adaptive accuracy gains coexist with slower recovery. Evaluating the joint outcome vector prevents either observation from being mistaken for uniform improvement.
